% =====================================================================
%  IEL04 - Circuitos Eléctricos I
%  Semana 12: El circuito RLC en paralelo: susceptancias opuestas, ley de corrientes de Kirchhoff, circuito tanque y corrección del factor de potencia
%  Institución Universitaria de Barranquilla (IUB)
%  Generada por src/presentaciones.py (diseño IUB 5). Compilador: pdfLaTeX (dos pasadas)
% =====================================================================
\documentclass[aspectratio=169,11pt,xcolor={table}]{beamer}

% ---------------------------------------------------------------------
%  Codificación, idioma y tipografía
% ---------------------------------------------------------------------
\usepackage[utf8]{inputenc}
\usepackage[T1]{fontenc}
\usepackage[spanish,es-nodecimaldot,es-noshorthands]{babel}
\usepackage{avant}                       % Sans geométrica (emula Avenir Next)
\renewcommand{\familydefault}{\sfdefault}
\renewcommand{\ttdefault}{pcr}           % Monoespaciada para código
\usepackage{textcomp}

% ---------------------------------------------------------------------
%  Paquetes
% ---------------------------------------------------------------------
\usepackage{amsmath,amssymb,mathtools}
\usepackage{siunitx}
\sisetup{output-decimal-marker={.}, per-mode=symbol}
\usepackage{booktabs,tabularx,multirow,array}
\usepackage{adjustbox}
\usepackage{tikz}
\usetikzlibrary{arrows.meta,positioning,calc,shapes.geometric,shapes.misc,
  decorations.pathreplacing,fit,backgrounds,babel}
\usepackage[siunitx,RPvoltages]{circuitikz}
\usepackage{pgfplots}
\pgfplotsset{compat=1.18}
\usepackage[most]{tcolorbox}
\usepackage{listings}

% ---------------------------------------------------------------------
%  Paleta institucional IUB (Manual de Identidad Visual 2025)
% ---------------------------------------------------------------------
\definecolor{iubwhite}{HTML}{FFFFFF}
\definecolor{iubnavy}{HTML}{121023}
\definecolor{iubyellow}{HTML}{FFDF2D}
\definecolor{iubblue}{HTML}{00ADE7}
\definecolor{iuborange}{HTML}{E39037}
\definecolor{iubgreen}{HTML}{3B9C6D}

% ---------------------------------------------------------------------
%  Configuración de Beamer
% ---------------------------------------------------------------------
\setbeamertemplate{navigation symbols}{}
\setbeamersize{text margin left=0.6cm, text margin right=0.6cm}
\setbeamercolor{normal text}{fg=iubnavy, bg=iubwhite}
\setbeamercolor{background canvas}{bg=iubwhite}
\setbeamercolor{structure}{fg=iubnavy}
\setbeamercolor{alerted text}{fg=iuborange}
\setbeamertemplate{itemize item}{\textcolor{iubblue}{\small$\blacktriangleright$}}
\setbeamertemplate{itemize subitem}{\textcolor{iuborange}{\scriptsize$\bullet$}}
\setbeamertemplate{enumerate item}{\textcolor{iubblue}{\bfseries\insertenumlabel.}}
\setbeamertemplate{frametitle}{}         % el título vive en el headline

% Parches: el headline se pre-mide antes del primer frame
\makeatletter
\providecommand{\insertframetitle}{}
\providecommand{\insertframesubtitle}{}
\makeatother

% ---------- Encabezado ----------
\setbeamertemplate{headline}{%
  \begin{tikzpicture}
    \useasboundingbox (0,0) rectangle (\paperwidth,1.05cm);
    \fill[iubnavy] (0,0) rectangle (\paperwidth,1.05cm);
    \fill[iubyellow] (0,0) rectangle (\paperwidth,0.05cm);
    \node[anchor=west, text=iubyellow, font=\large\bfseries]
      at (0.6cm,0.55cm) {\strut\insertframetitle};
    \node[anchor=east, text=iubyellow, font=\footnotesize\bfseries]
      at ([xshift=-0.6cm]\paperwidth,0.55cm) {IUB};
  \end{tikzpicture}}

% ---------- Pie de página ----------
\setbeamertemplate{footline}{%
  \begin{tikzpicture}
    \useasboundingbox (0,0) rectangle (\paperwidth,0.55cm);
    \fill[iubnavy] (0,0) rectangle (\paperwidth,0.55cm);
    \node[anchor=west, text=iubyellow, font=\tiny]
      at (0.6cm,0.275cm) {IEL04 \textperiodcentered{} Circuitos Eléctricos I};
    \node[anchor=center, text=iubyellow, font=\tiny]
      at (0.66\paperwidth,0.275cm) {Ing. Sergio Martinez Castilla};
    \node[anchor=east, text=iubyellow, font=\tiny\bfseries]
      at ([xshift=-0.6cm]\paperwidth,0.275cm)
      {Semana 12 \quad \insertframenumber\,/\,\inserttotalframenumber};
  \end{tikzpicture}}

% ---------------------------------------------------------------------
%  Cajas tcolorbox (texto blanco sobre la paleta complementaria)
%  \begin{caja}[opciones]{color}{título} ... \end{caja}
%  \begin{cajaplana}[opciones]{color} ... \end{cajaplana}
%  \begin{cardB|cardO|cardG|cardN}[opciones]{título} ... \end{cardX}
% ---------------------------------------------------------------------
\tcbset{
  iubbase/.style={enhanced, arc=1.5mm, boxrule=0pt, coltext=white, coltitle=white,
    fonttitle=\bfseries\footnotesize, fontupper=\footnotesize,
    left=1.8mm, right=1.8mm, top=1mm, bottom=1.2mm,
    toptitle=0.6mm, bottomtitle=0.6mm, halign=left,
    before upper={\hyphenpenalty=5000\setbeamercolor{itemize item}{fg=white}%
                  \setbeamercolor{itemize/enumerate body}{fg=white}%
                  \setbeamercolor{itemize/enumerate subbody}{fg=white}%
                  \setbeamercolor{itemize subitem}{fg=white}%
                  \setbeamercolor{enumerate item}{fg=white}%
                  \setbeamertemplate{itemize item}{\textcolor{white}{\small$\blacktriangleright$}}}}
}
\newtcolorbox{caja}[3][]{iubbase, colback=#2, colframe=#2,
  colbacktitle=#2!78!iubnavy, title={#3}, #1}
\newtcolorbox{cajaplana}[2][]{iubbase, colback=#2, colframe=#2, #1}
\newtcolorbox{cardB}[2][]{iubbase, colback=iubblue,   colframe=iubblue,
  colbacktitle=iubblue!78!iubnavy,   title={#2}, #1}
\newtcolorbox{cardO}[2][]{iubbase, colback=iuborange, colframe=iuborange,
  colbacktitle=iuborange!78!iubnavy, title={#2}, #1}
\newtcolorbox{cardG}[2][]{iubbase, colback=iubgreen,  colframe=iubgreen,
  colbacktitle=iubgreen!78!iubnavy,  title={#2}, #1}
\newtcolorbox{cardN}[2][]{iubbase, colback=iubnavy,   colframe=iubnavy,
  colbacktitle=iubnavy, coltitle=iubyellow, title={#2}, #1}

% Celda de encabezado de tabla (texto amarillo sobre \rowcolor{iubnavy}) y código en línea
\newcommand{\hd}[1]{\textcolor{iubyellow}{\bfseries #1}}
\newcommand{\thc}[1]{\textcolor{iubyellow}{\textbf{#1}}}
\newcommand{\cod}[1]{\texttt{\textbf{#1}}}

% ---------------------------------------------------------------------
%  Código sobre fondo oscuro:
%  \begin{codebox}[listing options app={language=Python,firstnumber=1}]{título}
% ---------------------------------------------------------------------
\lstdefinestyle{iubcode}{
  basicstyle=\ttfamily\fontsize{7}{8.4}\selectfont\color{white},
  keywordstyle=\color{iubblue}\bfseries,
  commentstyle=\color{iubgreen!55!white}\itshape,
  stringstyle=\color{iuborange!80!white},
  numbers=left, numberstyle=\tiny\color{iubyellow}, numbersep=6pt,
  showstringspaces=false, breaklines=true, tabsize=4,
  columns=fullflexible, keepspaces=true, upquote=true}
\newtcblisting{codebox}[2][]{enhanced, listing only, colback=iubnavy,
  colframe=iubnavy, colbacktitle=iubnavy!85!white, coltitle=iubyellow,
  fonttitle=\bfseries\scriptsize, title={#2}, arc=1.5mm, boxrule=0pt,
  left=6mm, right=1mm, top=0.5mm, bottom=0.5mm, toptitle=0.5mm, bottomtitle=0.5mm,
  listing options={style=iubcode}, #1}

% ---------------------------------------------------------------------
%  Estilos TikZ
% ---------------------------------------------------------------------
\tikzset{
  bloque/.style={rectangle, rounded corners=2mm, fill=#1, text=white,
    font=\scriptsize\bfseries, align=center, minimum height=1.1cm,
    text width=1.8cm, inner sep=1.2mm},
  flecha/.style={-{Stealth[length=2.2mm]}, line width=1pt, draw=iubnavy},
  arr/.style={flecha},
  term/.style={rectangle, draw=#1, line width=1.2pt, fill=white,
    text=iubnavy, font=\tiny\bfseries, minimum width=1.05cm,
    minimum height=0.5cm, inner sep=1pt},
  nodo/.style={rectangle, rounded corners=1mm, fill=#1, text=white,
    font=\scriptsize\bfseries, minimum size=0.75cm, align=center},
  cable/.style={line width=1.4pt, draw=#1},
  box/.style={rectangle, rounded corners=1.5mm, draw=none, text=white,
    font=\scriptsize\bfseries, align=center, minimum height=0.8cm, inner sep=3pt},
  bB/.style={box, fill=iubblue}, bO/.style={box, fill=iuborange},
  bG/.style={box, fill=iubgreen}, bN/.style={box, fill=iubnavy},
  dec/.style={diamond, aspect=1.9, fill=iuborange, text=white,
    font=\scriptsize\bfseries, align=center, inner sep=1pt},
  tag/.style={fill=#1, text=white, font=\scriptsize\bfseries, rounded corners=1mm,
    minimum width=0.7cm, anchor=north west},
}

% ---------------------------------------------------------------------
%  Diapositiva separadora de bloque
%  #1 número  #2 título  #3 duración  #4 color
% ---------------------------------------------------------------------
\newcommand{\bloque}[4]{%
\section{#2}
\begin{frame}[plain]
\begin{tikzpicture}[remember picture, overlay]
  \fill[#4] (current page.north west) rectangle
        ([xshift=5.2cm]current page.south west);
  \node[anchor=north west, text=white, font=\bfseries\large]
    at ([xshift=0.8cm,yshift=-2.2cm]current page.north west) {BLOQUE};
  \node[anchor=north west, text=white, font=\bfseries\fontsize{72}{72}\selectfont]
    at ([xshift=0.65cm,yshift=-2.9cm]current page.north west) {#1};
  \node[anchor=north west, text=iubnavy, font=\bfseries\LARGE,
        text width=9.4cm, align=left]
    at ([xshift=6.2cm,yshift=-2.8cm]current page.north west)
    {\hyphenpenalty=10000\exhyphenpenalty=10000 #2};
  \node[anchor=north west, fill=iubnavy, text=iubyellow, rounded corners=1.5mm,
        font=\small\bfseries, inner sep=2mm]
    at ([xshift=6.2cm,yshift=-6.0cm]current page.north west) {Duración: #3};
  \fill[iubnavy] ([xshift=5.2cm]current page.south west) rectangle
        ([yshift=0.35cm]current page.south east);
  \fill[iubyellow] ([xshift=5.2cm,yshift=0.35cm]current page.south west)
        rectangle ([yshift=0.42cm]current page.south east);
  \node[anchor=north east, text=iubnavy, font=\small\bfseries]
    at ([xshift=-0.6cm,yshift=-0.5cm]current page.north east) {IUB · IEL04};
\end{tikzpicture}
\end{frame}}

% =====================================================================
\begin{document}
% =====================================================================

% ---------------------------------------------------------------------
%  PORTADA
% ---------------------------------------------------------------------
\begin{frame}[plain]
\begin{tikzpicture}[remember picture, overlay]
  % Panel institucional
  \fill[iubnavy] (current page.north west) rectangle
        ([xshift=5.6cm]current page.south west);
  \node[anchor=north west, text=iubyellow, font=\bfseries\fontsize{54}{54}\selectfont]
    at ([xshift=0.7cm,yshift=-1.0cm]current page.north west) {IUB};
  \node[anchor=north west, text=white, font=\footnotesize, text width=4.3cm, align=left]
    at ([xshift=0.75cm,yshift=-3.0cm]current page.north west)
    {Institución Universitaria de Barranquilla};
  \draw[iubyellow, line width=1.2pt]
    ([xshift=0.75cm,yshift=-4.25cm]current page.north west) --
    ([xshift=2.6cm,yshift=-4.25cm]current page.north west);
  \node[anchor=north west, text=iubyellow, font=\scriptsize\bfseries,
        text width=4.3cm, align=left]
    at ([xshift=0.75cm,yshift=-4.5cm]current page.north west)
    {Tecnología en Gestión de Sistemas Eléctricos};
  \node[anchor=south west, text=white, font=\scriptsize, text width=4.3cm, align=left]
    at ([xshift=0.75cm,yshift=0.9cm]current page.south west)
    {Ing. Sergio Martinez Castilla\\[1pt]\textcolor{iubyellow}{\tiny smartinezc@unibarranquilla.edu.co}};
  % Bloque de título
  \node[anchor=north west, fill=iubblue, text=white, rounded corners=1.5mm,
        font=\scriptsize\bfseries, inner sep=2mm]
    at ([xshift=6.4cm,yshift=-1.0cm]current page.north west)
    {IEL04 \textperiodcentered{} Semana 12 de 14 \textperiodcentered{} Corte evaluativo 3};
  \node[anchor=north west, text=iubnavy, font=\small, text width=9.0cm, align=left] (a)
    at ([xshift=6.4cm,yshift=-1.9cm]current page.north west)
    {Circuitos Eléctricos I};
  \node[anchor=north west, text=iubnavy, font=\bfseries\fontsize{13}{16}\selectfont,
        text width=9.0cm, align=left] (t)
    at ([yshift=-0.35cm]a.south west)
    {\hyphenpenalty=10000 El circuito RLC en paralelo: susceptancias opuestas, ley de corrientes de Kirchhoff, circuito tanque y corrección del factor de potencia};
  % Franjas de la paleta complementaria
  \fill[iubblue]   ([xshift=6.4cm,yshift=1.1cm]current page.south west)
    rectangle ++(1.6cm,0.18cm);
  \fill[iuborange] ([xshift=8.1cm,yshift=1.1cm]current page.south west)
    rectangle ++(1.6cm,0.18cm);
  \fill[iubgreen]  ([xshift=9.8cm,yshift=1.1cm]current page.south west)
    rectangle ++(1.6cm,0.18cm);
  \node[anchor=north west, text=iubnavy, font=\scriptsize]
    at ([xshift=6.4cm,yshift=0.95cm]current page.south west)
    {Sesión teórico-práctica \textperiodcentered{} 240 min};
\end{tikzpicture}
\end{frame}

% ---------------------------------------------------------------------
\begin{frame}{Punto de partida de la sesión}
\begin{columns}[T,onlytextwidth]
  \column{0.34\textwidth}
\begin{caja}[height=5.9cm]{iubblue}{Continuidad}
\scriptsize \begin{itemize}\setlength{\itemsep}{2pt}
  \item Semana 11: Circuito RLC en serie: análisis y Ley de Voltajes de Kirchhoff
  \item \textbf{Semana 12: Circuito RLC en paralelo: análisis y Ley de Corrientes de Kirchhoff}
  \item Semana 13: Cálculo y medición de parámetros en circuitos RLC serie y paralelo
\end{itemize}
\end{caja}
  \column{0.34\textwidth}
\begin{caja}[height=5.9cm]{iuborange}{Prerrequisitos}
\scriptsize \begin{itemize}\setlength{\itemsep}{2pt}
  \item Calcular admitancias y aplicar la ley de corrientes de Kirchhoff
  \item Analizar un circuito RLC en serie y su resonancia (semana 11)
  \item Convertir una bobina real a su modelo en paralelo con R\_p
  \item Calcular potencia activa, reactiva y aparente y factor de potencia
\end{itemize}
\end{caja}
  \column{0.29\textwidth}
\begin{caja}[height=5.9cm]{iubgreen}{Competencia}
\scriptsize \textbf{UC2.} Coordinar actividades asociadas a sistemas eléctricos utilizando fuentes convencionales y no convencionales de energía.
\end{caja}
\end{columns}
\end{frame}

% ---------------------------------------------------------------------
\begin{frame}{Resultados de aprendizaje}
\begin{columns}[c,onlytextwidth]
  \column{0.45\textwidth}
\begin{caja}{iubnavy}{IEL04-RA4}
\footnotesize Calcular un circuito resistivo, inductivo, capacitivo y RLC, sea en serie o paralelo.
\end{caja}
\vspace{1mm}
\begin{cajaplana}{iubblue}
\scriptsize \textbf{CE1.} Identificar los tipos de resistores, inductores y capacitores.\\[2pt]
\textbf{CE2.} Realizar mediciones en un circuito resistivo, inductivo y capacitivo.\\[2pt]
\textbf{CE3.} Realizar mediciones en un circuito RLC.
\end{cajaplana}
  \column{0.52\textwidth}
\centering
\begin{adjustbox}{max width=\linewidth, max totalheight=6.1cm}
\begin{tikzpicture}
  \node[box, fill=iubblue, text width=4.9cm, align=left, font=\tiny, anchor=west] (r0) at (1.25,0.00) {Identificar y verificar el resistor, la bobina y el condensador de un circuito RLC en paralelo antes del montaje.};
  \node[tag=iubnavy, font=\tiny\bfseries, minimum width=1.15cm, anchor=east] at (1.1,0.00) {Aplicar};
  \node[box, fill=iuborange, text width=4.9cm, align=left, font=\tiny, anchor=west] (r1) at (1.25,-1.45) {Calcular la admitancia, la impedancia y el carácter de un circuito RLC en paralelo a una frecuencia dada.};
  \node[tag=iubnavy, font=\tiny\bfseries, minimum width=1.15cm, anchor=east] at (1.1,-1.45) {Aplicar};
  \node[box, fill=iubgreen, text width=4.9cm, align=left, font=\tiny, anchor=west] (r2) at (1.25,-2.90) {Analizar las corrientes de un circuito RLC en paralelo con la ley de corrientes de Kirchhoff y explicar la resonancia en paralelo y el circuito tanque real.};
  \node[tag=iubnavy, font=\tiny\bfseries, minimum width=1.15cm, anchor=east] at (1.1,-2.90) {Analizar};
  \node[box, fill=iubblue, text width=4.9cm, align=left, font=\tiny, anchor=west] (r3) at (1.25,-4.35) {Medir las corrientes de rama y la corriente total de un RLC en paralelo a varias frecuencias y verificar la resonancia.};
  \node[tag=iubnavy, font=\tiny\bfseries, minimum width=1.15cm, anchor=east] at (1.1,-4.35) {Evaluar};
  \draw[flecha, draw=iubnavy!40] (r0.south) -- (r1.north);
  \draw[flecha, draw=iubnavy!40] (r1.south) -- (r2.north);
  \draw[flecha, draw=iubnavy!40] (r2.south) -- (r3.north);
\end{tikzpicture}
\end{adjustbox}
\end{columns}
\end{frame}

% ---------------------------------------------------------------------
\begin{frame}{Ruta de la sesión \textperiodcentered{} 240 minutos}
\centering
\begin{tikzpicture}[x=1cm,y=1cm]
  \node[circle, fill=iubblue, text=white, font=\scriptsize\bfseries, minimum size=0.6cm, inner sep=0] at (0,0.00) {1};
  \node[anchor=west, font=\scriptsize, text width=7.6cm] at (0.45,0.00) {Del RLC en serie al RLC en paralelo: tensión común y corrientes reactivas opuestas};
  \fill[iubblue] (8.3,-0.20) rectangle ++(2.04,0.4);
  \node[anchor=west, font=\scriptsize\bfseries] at (10.44,0.00) {20 min};
  \node[circle, fill=iuborange, text=white, font=\scriptsize\bfseries, minimum size=0.6cm, inner sep=0] at (0,-0.76) {2};
  \node[anchor=west, font=\scriptsize, text width=7.6cm] at (0.45,-0.76) {Y = G + j(BC \ensuremath{-} BL), carácter inductivo o capacitivo según la frecuencia};
  \fill[iuborange] (8.3,-0.96) rectangle ++(3.07,0.4);
  \node[anchor=west, font=\scriptsize\bfseries] at (11.47,-0.76) {30 min};
  \node[circle, fill=iubgreen, text=white, font=\scriptsize\bfseries, minimum size=0.6cm, inner sep=0] at (0,-1.51) {3};
  \node[anchor=west, font=\scriptsize, text width=7.6cm] at (0.45,-1.51) {IT = IR + j(IC \ensuremath{-} IL): corrientes reactivas en oposición y mayores que la total};
  \fill[iubgreen] (8.3,-1.71) rectangle ++(3.58,0.4);
  \node[anchor=west, font=\scriptsize\bfseries] at (11.98,-1.51) {35 min};
  \node[circle, fill=iubblue, text=white, font=\scriptsize\bfseries, minimum size=0.6cm, inner sep=0] at (0,-2.27) {4};
  \node[anchor=west, font=\scriptsize, text width=7.6cm] at (0.45,-2.27) {fr, Z = R, Q = R/XL, ancho de banda y curva de impedancia};
  \fill[iubblue] (8.3,-2.47) rectangle ++(3.58,0.4);
  \node[anchor=west, font=\scriptsize\bfseries] at (11.98,-2.27) {35 min};
  \node[circle, fill=iuborange, text=white, font=\scriptsize\bfseries, minimum size=0.6cm, inner sep=0] at (0,-3.03) {5};
  \node[anchor=west, font=\scriptsize, text width=7.6cm] at (0.45,-3.03) {Equivalente RLC en paralelo, Rp(eq) = RW(Q\texttwosuperior{} + 1) y desplazamiento de la resonancia};
  \fill[iuborange] (8.3,-3.23) rectangle ++(3.07,0.4);
  \node[anchor=west, font=\scriptsize\bfseries] at (11.47,-3.03) {30 min};
  \node[circle, fill=iubgreen, text=white, font=\scriptsize\bfseries, minimum size=0.6cm, inner sep=0] at (0,-3.79) {6};
  \node[anchor=west, font=\scriptsize, text width=7.6cm] at (0.45,-3.79) {Compensar la corriente reactiva de una carga inductiva y dimensionar el condensador};
  \fill[iubgreen] (8.3,-3.99) rectangle ++(3.07,0.4);
  \node[anchor=west, font=\scriptsize\bfseries] at (11.47,-3.79) {30 min};
  \node[circle, fill=iubblue, text=white, font=\scriptsize\bfseries, minimum size=0.6cm, inner sep=0] at (0,-4.54) {7};
  \node[anchor=west, font=\scriptsize, text width=7.6cm] at (0.45,-4.54) {Con el kit: 4.7 k\ensuremath{\Omega}, bobina de 10.4 mH y condensador de 97.8 nF a 3 kHz, fr y 8 kHz};
  \fill[iubblue] (8.3,-4.74) rectangle ++(4.60,0.4);
  \node[anchor=west, font=\scriptsize\bfseries] at (13.00,-4.54) {45 min};
  \node[circle, fill=iuborange, text=white, font=\scriptsize\bfseries, minimum size=0.6cm, inner sep=0] at (0,-5.30) {8};
  \node[anchor=west, font=\scriptsize, text width=7.6cm] at (0.45,-5.30) {Síntesis, cierre actitudinal y bibliografía};
  \fill[iuborange] (8.3,-5.50) rectangle ++(1.53,0.4);
  \node[anchor=west, font=\scriptsize\bfseries] at (9.93,-5.30) {15 min};
  \draw[iubnavy, line width=0.6pt] (8.3,0.45) -- (8.3,-5.70);
\end{tikzpicture}
\end{frame}

% =====================================================================
\bloque{01}{Del RLC en serie al RLC en paralelo: tensión común y corrientes reactivas opuestas}{20 min}{iubblue}
% =====================================================================

% ---------------------------------------------------------------------
\begin{frame}{{\normalsize Tres elementos en paralelo: cuando las corrientes se oponen}}
\begin{columns}[c,onlytextwidth]
  \column{0.57\textwidth}
\small
\begin{itemize}\setlength{\itemsep}{4pt}
  \item En la semana 11, el resistor, la bobina y el condensador compartían la corriente, y sus tensiones reactivas, opuestas, se restaban.
  \item El circuito RLC en paralelo aparece en tres contextos que el técnico en sistemas eléctricos encuentra a menudo.
  \item Cuando son iguales, en la resonancia, se cancelan y la fuente solo entrega la corriente del resistor: la corriente total es mínima y la impedancia máxima.
\end{itemize}
  \column{0.4\textwidth}
\begin{cardO}{Nota pedagógica}
\footnotesize Conviene tener a la vista la tabla de la semana 11 al estudiar esta sesión: casi cada resultado del RLC en paralelo es el del RLC en serie con tensión y corriente, impedancia y admitancia, y mínimo y máximo intercambiados.
\end{cardO}
\end{columns}
\end{frame}

% ---------------------------------------------------------------------
\begin{frame}{Relaciones de fase en el circuito RLC en paralelo}
\centering
\begin{adjustbox}{max width=\linewidth, max totalheight=6.1cm}
\begin{tikzpicture}
  \node[bG, align=center, font=\scriptsize, text width=2.10cm] (n0) at (1.17,2.80) {Tensión común VS};
  \node[bB, align=center, font=\scriptsize, text width=2.07cm] (n1) at (4.65,2.80) {IR en fase con VS};
  \node[bB, align=center, font=\scriptsize, text width=1.90cm] (n2) at (4.65,0.81) {IL retrasada 90°};
  \node[bB, align=center, font=\scriptsize, text width=2.07cm] (n3) at (4.65,4.79) {IC adelantada 90°};
  \node[bO, align=center, font=\scriptsize, text width=2.24cm] (n4) at (8.19,2.80) {LCK: IT = IR + j(IC \ensuremath{-} IL)};
  \node[bG, align=center, font=\scriptsize, text width=2.44cm] (n5) at (11.92,2.80) {IT mínima en resonancia};
  \draw[flecha] (n0) -- (n1);
  \draw[flecha] (n0) -- (n2);
  \draw[flecha] (n0) -- (n3);
  \draw[flecha] (n1) -- (n4);
  \draw[flecha] (n2) -- (n4);
  \draw[flecha] (n3) -- (n4);
  \draw[flecha] (n4) -- (n5);
\end{tikzpicture}
\end{adjustbox}

\vspace{1mm}{\scriptsize Relaciones de fase en el circuito RLC en paralelo: de la tensión común a la corriente total\par}
\end{frame}

% =====================================================================
\bloque{02}{Y = G + j(BC \ensuremath{-} BL), carácter inductivo o capacitivo según la frecuencia}{30 min}{iuborange}
% =====================================================================

% ---------------------------------------------------------------------
\begin{frame}{{\normalsize Admitancia e impedancia del circuito RLC en paralelo}}
\begin{columns}[c,onlytextwidth]
  \column{0.57\textwidth}
\small
\begin{itemize}\setlength{\itemsep}{4pt}
  \item En paralelo, como se vio en las semanas 6 y 8, conviene trabajar con admitancias: se suman directamente, igual que las impedancias en serie.
  \item Lejos de la resonancia, una de las dos ramas reactivas deriva casi toda la corriente y la impedancia cae a unos cientos de ohmios.
  \item Si el circuito lleva una bobina real, su resistencia de devanado no está en paralelo sino en serie con la inductancia, dentro de la misma rama.
\end{itemize}
  \column{0.4\textwidth}
\begin{caja}{iubblue}{Ecuación clave}
\footnotesize \centering\small \begin{adjustbox}{max width=\linewidth}$\displaystyle \begin{gathered}\mathbf{Y}= G + j(B_C - B_L) \\ = \sqrt{G^{2} + (B_C - B_L)^{2}}\ \angle \tan^{-1}\!\left(\frac{B_C - B_L}{G}\right) \\ \mathbf{Z} = \frac{1}{\mathbf{Y}}\end{gathered}$\end{adjustbox}
\end{caja}
\end{columns}
\end{frame}

% ---------------------------------------------------------------------
\begin{frame}{{\normalsize Admitancia e impedancia de un RLC en paralelo ideal}}
\centering\tiny
\renewcommand{\arraystretch}{1.35}
\rowcolors{2}{iubnavy!6}{white}
\begin{adjustbox}{max width=\linewidth, max totalheight=5.6cm}
\begin{tabularx}{\textwidth}{>{\raggedright\arraybackslash}X>{\raggedright\arraybackslash}X>{\raggedright\arraybackslash}X>{\raggedright\arraybackslash}X>{\raggedright\arraybackslash}X>{\raggedright\arraybackslash}X>{\raggedright\arraybackslash}X>{\raggedright\arraybackslash}X}
  \rowcolor{iubnavy}
  \hd{f (Hz)} & \hd{G (mS)} & \hd{BC (mS)} & \hd{BL (mS)} & \hd{Y (mS)} & \hd{Z (\ensuremath{\Omega})} & \hd{Ángulo de Z (°)} & \hd{Carácter}\\
  2000 & 1.00 & 1.26 & 7.96 & 6.78 & 147.6 & +81.5 & inductivo\\
  4000 & 1.00 & 2.51 & 3.98 & 1.78 & 563.5 & +55.7 & inductivo\\
  5033 (fr) & 1.00 & 3.16 & 3.16 & 1.00 & 1000 & 0.0 & resistivo\\
  6000 & 1.00 & 3.77 & 2.65 & 1.50 & 666.9 & \ensuremath{-}48.2 & capacitivo\\
  10 000 & 1.00 & 6.28 & 1.59 & 4.80 & 208.5 & \ensuremath{-}78.0 & capacitivo\\
\end{tabularx}
\end{adjustbox}
\vspace{2mm}
{\scriptsize Admitancia e impedancia de un RLC en paralelo ideal (R = 1 k\ensuremath{\Omega}, L = 10 mH, C = 0.1 µF) a varias frecuencias\par}
\end{frame}

% =====================================================================
\bloque{03}{IT = IR + j(IC \ensuremath{-} IL): corrientes reactivas en oposición y mayores que la total}{35 min}{iubgreen}
% =====================================================================

% ---------------------------------------------------------------------
\begin{frame}{Ley de corrientes de Kirchhoff en el circuito RLC}
\begin{columns}[c,onlytextwidth]
  \column{0.57\textwidth}
\small
\begin{itemize}\setlength{\itemsep}{4pt}
  \item La ley de corrientes de Kirchhoff se aplica en el nodo donde la corriente de la fuente se divide entre las tres ramas.
  \item Con $R = 1\ \text{k}\Omega$, $L = 10\ \text{mH}$, $C = 0.1\ \mu\text{F}$ y 2 V eficaces, a 4 kHz: $I_R = 2.00\ \text{mA}$, $I_C = 2 \times 2.513 = 5.03\ \text{mA}$ e $I_L = 2 \times 3.979 = 7.96\ \text{mA}$.
  \item Las columnas de la tabla también muestran por qué el carácter del circuito cambia con la frecuencia.
\end{itemize}
  \column{0.4\textwidth}
\begin{caja}{iubblue}{Ecuación clave}
\footnotesize \centering\small \begin{adjustbox}{max width=\linewidth}$\displaystyle \begin{gathered}\mathbf{I}_T = I_R + j(I_C - I_L) \\ I_T = \sqrt{I_R^{2} + (I_C - I_L)^{2}} \\ \theta = \tan^{-1}\!\left(\frac{I_C - I_L}{I_R}\right)\end{gathered}$\end{adjustbox}
\end{caja}
\end{columns}
\end{frame}

% ---------------------------------------------------------------------
\begin{frame}{Corrientes de un RLC en paralelo ideal}
\centering\tiny
\renewcommand{\arraystretch}{1.35}
\rowcolors{2}{iubnavy!6}{white}
\begin{adjustbox}{max width=\linewidth, max totalheight=5.6cm}
\begin{tabularx}{\textwidth}{>{\raggedright\arraybackslash}X>{\raggedright\arraybackslash}X>{\raggedright\arraybackslash}X>{\raggedright\arraybackslash}X>{\raggedright\arraybackslash}X>{\raggedright\arraybackslash}X>{\raggedright\arraybackslash}X}
  \rowcolor{iubnavy}
  \hd{f (Hz)} & \hd{IR (mA)} & \hd{IL (mA)} & \hd{IC (mA)} & \hd{IT (mA)} & \hd{\ensuremath{\theta} (°)} & \hd{Carácter}\\
  2000 & 2.00 & 15.92 & 2.51 & 13.56 & \ensuremath{-}81.5 & inductivo\\
  4000 & 2.00 & 7.96 & 5.03 & 3.55 & \ensuremath{-}55.7 & inductivo\\
  5033 (fr) & 2.00 & 6.32 & 6.32 & 2.00 & 0.0 & resistivo\\
  10 000 & 2.00 & 3.18 & 12.57 & 9.59 & +78.0 & capacitivo\\
\end{tabularx}
\end{adjustbox}
\vspace{2mm}
{\scriptsize Corrientes de un RLC en paralelo ideal (R = 1 k\ensuremath{\Omega}, L = 10 mH, C = 0.1 µF, VS = 2 V eficaces)\par}
\end{frame}

% =====================================================================
\bloque{04}{fr, Z = R, Q = R/XL, ancho de banda y curva de impedancia}{35 min}{iubblue}
% =====================================================================

% ---------------------------------------------------------------------
\begin{frame}{{\normalsize Resonancia en paralelo: impedancia máxima y corriente mínima}}
\begin{columns}[c,onlytextwidth]
  \column{0.57\textwidth}
\small
\begin{itemize}\setlength{\itemsep}{4pt}
  \item La resonancia en paralelo ocurre cuando las susceptancias capacitiva e inductiva son iguales, $B_C = B_L$, lo que equivale a $X_L = X_C$.
  \item Con $L = 10\ \text{mH}$ y $C = 0.1\ \mu\text{F}$, $f_r = 5.03\ \text{kHz}$ y $X_L = 316\ \Omega$.
\end{itemize}
  \column{0.4\textwidth}
\begin{caja}{iubblue}{Ecuación clave}
\footnotesize \centering\small \begin{adjustbox}{max width=\linewidth}$\displaystyle \begin{gathered}f_r = \frac{1}{2\pi\sqrt{LC}} \\ Z_r = R \\ I_{T,min} = \frac{V_S}{R} \\ Q = \frac{R}{X_L} \\ AB = \frac{f_r}{Q}\end{gathered}$\end{adjustbox}
\end{caja}
\end{columns}
\end{frame}

% ---------------------------------------------------------------------
\begin{frame}{{\normalsize Impedancia de un RLC en paralelo (L = 10 mH, C = 0.1 µF)}}
\begin{columns}[c,onlytextwidth]
  \column{0.62\textwidth}
\centering
\begin{tikzpicture}
  \begin{semilogxaxis}[width=8.4cm, height=5.7cm, xmin=1000, xmax=25000, ymin=0, ymax=2962.6,
      xlabel={Frecuencia f (Hz)}, ylabel={Impedancia Z (\ensuremath{\Omega})},
      label style={font=\scriptsize, text=iubnavy}, tick label style={font=\tiny, text=iubnavy},
      axis line style={iubnavy}, grid=major, grid style={iubnavy!12},
      legend style={font=\tiny, fill=white}, legend pos=north east]
    \addplot[line width=1.4pt, iubblue] coordinates {(1000,65.275) (1044,68.386) (1089.9,71.668) (1137.9,75.135) (1188,78.801) (1240.2,82.681) (1294.8,86.794) (1351.8,91.158) (1411.3,95.798) (1473.4,100.74) (1538.2,106.01) (1605.9,111.64) (1676.6,117.67) (1750.3,124.14) (1827.3,131.11) (1907.8,138.63) (1991.7,146.77) (2079.3,155.62) (2170.8,165.27) (2266.4,175.84) (2366.1,187.46) (2470.2,200.32) (2578.9,214.61) (2692.4,230.6) (2810.8,248.61) (2934.5,269.05) (3063.7,292.44) (3198.5,319.45) (3339.2,350.95) (3486.1,388.06) (3639.5,432.26) (3799.7,485.39) (3966.9,549.71) (4141.4,627.49) (4323.7,719.81) (4513.9,823.2) (4712.5,923.2) (4919.9,989.84) (5136.4,991.82) (5362.4,928.05) (5598.4,828.97) (5844.7,725.26) (6101.9,632.18) (6370.4,553.61) (6650.7,488.6) (6943.4,434.92) (7248.9,390.29) (7567.8,352.83) (7900.9,321.05) (8248.5,293.82) (8611.5,270.25) (8990.4,249.66) (9386,231.53) (9799,215.44) (10230,201.06) (10680,188.13) (11150,176.44) (11641,165.82) (12153,156.13) (12688,147.24) (13246,139.06) (13829,131.51) (14438,124.51) (15073,118.01) (15736,111.96) (16429,106.3) (17151,101.02) (17906,96.061) (18694,91.405) (19517,87.026) (20375,82.9) (21272,79.008) (22208,75.331) (23185,71.853) (24205,68.561) (25000,66.206)};
    \addlegendentry{R = 1 k\ensuremath{\Omega} (Q = 3.16)}
    \addplot[line width=1.4pt, iuborange] coordinates {(1000,65.399) (1044,68.528) (1089.9,71.832) (1137.9,75.324) (1188,79.019) (1240.2,82.933) (1294.8,87.086) (1351.8,91.497) (1411.3,96.191) (1473.4,101.19) (1538.2,106.54) (1605.9,112.26) (1676.6,118.4) (1750.3,125) (1827.3,132.12) (1907.8,139.83) (1991.7,148.2) (2079.3,157.32) (2170.8,167.31) (2266.4,178.3) (2366.1,190.46) (2470.2,203.99) (2578.9,219.14) (2692.4,236.25) (2810.8,255.73) (2934.5,278.15) (3063.7,304.24) (3198.5,335.01) (3339.2,371.9) (3486.1,416.97) (3639.5,473.35) (3799.7,545.91) (3966.9,642.77) (4141.4,778.31) (4323.7,980.05) (4513.9,1305.5) (4712.5,1875.1) (4919.9,2755.1) (5136.4,2798.6) (5362.4,1916.8) (5598.4,1328.8) (5844.7,993.93) (6101.9,787.3) (6370.4,649.03) (6650.7,550.49) (6943.4,476.85) (7248.9,419.73) (7567.8,374.13) (7900.9,336.86) (8248.5,305.79) (8611.5,279.48) (8990.4,256.88) (9386,237.25) (9799,220.03) (10230,204.77) (10680,191.16) (11150,178.94) (11641,167.89) (12153,157.85) (12688,148.68) (13246,140.27) (13829,132.53) (14438,125.38) (15073,118.75) (15736,112.59) (16429,106.84) (17151,101.48) (17906,96.457) (18694,91.747) (19517,87.32) (20375,83.155) (21272,79.228) (22208,75.521) (23185,72.019) (24205,68.705) (25000,66.335)};
    \addlegendentry{R = 3 k\ensuremath{\Omega} (Q = 9.49)}
    \addplot[line width=1pt, dashed, iubnavy!70] coordinates {(5033,0) (5033,2962.6)};
    \addlegendentry{fr = 5.03 kHz}
    \addplot[line width=1pt, dashed, iubnavy!70] coordinates {(1000,707) (25000,707)};
    \addlegendentry{0.707 · Zmax (R = 1 k\ensuremath{\Omega})}
  \end{semilogxaxis}
\end{tikzpicture}
  \column{0.35\textwidth}
\begin{caja}{iubblue}{Lectura de la gráfica}
\footnotesize Y, a la inversa, un tanque en serie con una línea bloquea precisamente su frecuencia de resonancia, que es el principio de los filtros de rechazo de banda.
\end{caja}
\end{columns}
\end{frame}

% =====================================================================
\bloque{05}{Equivalente RLC en paralelo, Rp(eq) = RW(Q\texttwosuperior{} + 1) y desplazamiento de la resonancia}{30 min}{iuborange}
% =====================================================================

% ---------------------------------------------------------------------
\begin{frame}{El circuito tanque con bobina real}
\begin{columns}[c,onlytextwidth]
  \column{0.57\textwidth}
\small
\begin{itemize}\setlength{\itemsep}{4pt}
  \item El circuito tanque ideal supone una bobina sin resistencia.
  \item La resistencia de devanado tiene además un efecto sobre la propia frecuencia de resonancia.
\end{itemize}
  \column{0.4\textwidth}
\begin{caja}{iubblue}{Ecuación clave}
\footnotesize \centering\small \begin{adjustbox}{max width=\linewidth}$\displaystyle \begin{gathered}Q = \frac{X_L}{R_W} \\ L_{eq} = L\left(\frac{Q^{2} + 1}{Q^{2}}\right) \\ R_{p(eq)} = R_W\,(Q^{2} + 1)\end{gathered}$\end{adjustbox}
\end{caja}
\end{columns}
\end{frame}

% ---------------------------------------------------------------------
\begin{frame}{{\normalsize Efecto de la resistencia de devanado en un tanque con L}}
\centering\scriptsize
\renewcommand{\arraystretch}{1.35}
\rowcolors{2}{iubnavy!6}{white}
\begin{adjustbox}{max width=\linewidth, max totalheight=5.6cm}
\begin{tabularx}{\textwidth}{>{\raggedright\arraybackslash}X>{\raggedright\arraybackslash}X>{\raggedright\arraybackslash}X>{\raggedright\arraybackslash}X>{\raggedright\arraybackslash}X}
  \rowcolor{iubnavy}
  \hd{RW (\ensuremath{\Omega})} & \hd{Q = XL/RW} & \hd{Rp(eq) = RW(Q² + 1) (\ensuremath{\Omega})} & \hd{fr con corriente en fase (Hz)} & \hd{Desplazamiento}\\
  24.6 (L3 del kit) & 13.3 & 4349 & 4976 & \ensuremath{-}0.3 \%\\
  50 & 6.52 & 2176 & 4931 & \ensuremath{-}1.2 \%\\
  100 & 3.26 & 1163 & 4750 & \ensuremath{-}4.8 \%\\
\end{tabularx}
\end{adjustbox}
\vspace{2mm}
{\scriptsize Efecto de la resistencia de devanado en un tanque con L = 10.4 mH y C = 97.8 nF (resonancia ideal: 4990 Hz, XL = 326.1 \ensuremath{\Omega})\par}
\end{frame}

% =====================================================================
\bloque{06}{Compensar la corriente reactiva de una carga inductiva y dimensionar el condensador}{30 min}{iubgreen}
% =====================================================================

% ---------------------------------------------------------------------
\begin{frame}{{\normalsize Corrección del factor de potencia con un condensador}}
\begin{columns}[c,onlytextwidth]
  \column{0.57\textwidth}
\small
\begin{itemize}\setlength{\itemsep}{4pt}
  \item La aplicación más importante del RLC en paralelo en los sistemas eléctricos es la \textbf{corrección del factor de potencia}.
  \item El condensador no cambia la potencia real de la carga, porque no disipa energía; solo reduce la potencia reactiva que la red debe suministrar.
  \item El ángulo inicial es $\cos^{-1}0.70 = 45.6^\circ$, con $\tan\theta_1 = 1.020$; el final, $\cos^{-1}0.95 = 18.2^\circ$, con $\tan\theta_2 = 0.329$.
\end{itemize}
  \column{0.4\textwidth}
\begin{caja}{iubblue}{Ecuación clave}
\footnotesize \centering\small \begin{adjustbox}{max width=\linewidth}$\displaystyle \begin{gathered}Q_C = P\,(\tan\theta_1 - \tan\theta_2) \\ C = \frac{Q_C}{2\pi f V^{2}}\end{gathered}$\end{adjustbox}
\end{caja}
\end{columns}
\end{frame}

% ---------------------------------------------------------------------
\begin{frame}{{\normalsize Corrección del factor de potencia de una carga de 1.2 kW}}
\centering\scriptsize
\renewcommand{\arraystretch}{1.35}
\rowcolors{2}{iubnavy!6}{white}
\begin{adjustbox}{max width=\linewidth, max totalheight=5.6cm}
\begin{tabularx}{\textwidth}{>{\raggedright\arraybackslash}X>{\raggedright\arraybackslash}X>{\raggedright\arraybackslash}X>{\raggedright\arraybackslash}X}
  \rowcolor{iubnavy}
  \hd{Magnitud} & \hd{Sin condensador} & \hd{Con 153 µF en paralelo} & \hd{Cambio}\\
  Potencia real P & 1200 W & 1200 W & sin cambio\\
  Potencia reactiva Q & 1224 VAR & 394 VAR & \ensuremath{-}68 \%\\
  Potencia aparente S & 1714 VA & 1263 VA & \ensuremath{-}26 \%\\
  Factor de potencia & 0.70 en atraso & 0.95 en atraso & mejora\\
  Corriente de línea & 14.3 A & 10.5 A & \ensuremath{-}26 \%\\
  Corriente del condensador & — & 6.9 A & —\\
\end{tabularx}
\end{adjustbox}
\vspace{2mm}
{\scriptsize Corrección del factor de potencia de una carga de 1.2 kW a 120 V y 60 Hz\par}
\end{frame}

% =====================================================================
\bloque{07}{Con el kit: 4.7 k\ensuremath{\Omega}, bobina de 10.4 mH y condensador de 97.8 nF a 3 kHz, fr y 8 kHz}{45 min}{iubblue}
% =====================================================================

% ---------------------------------------------------------------------
\begin{frame}{{\normalsize Medición de un circuito RLC en paralelo antes, en y después}}
\begin{columns}[c,onlytextwidth]
  \column{0.57\textwidth}
\small
\begin{itemize}\setlength{\itemsep}{4pt}
  \item Los tres se verifican antes del montaje con el ohmímetro y el medidor LCR, lo que cubre el criterio CE1.
  \item Con la bobina real, la corriente total se calcula sumando las admitancias de las tres ramas en forma rectangular.
  \item La comparación con el modelo ideal es instructiva.
\end{itemize}
  \column{0.4\textwidth}
\begin{caja}{iubblue}{Ecuación clave}
\footnotesize \centering\small \begin{adjustbox}{max width=\linewidth}$\displaystyle I_T = V_S\sqrt{\left(\frac{1}{R} + \frac{R_W}{R_W^{2} + X_L^{2}}\right)^{2} + \left(2\pi f C - \frac{X_L}{R_W^{2} + X_L^{2}}\right)^{2}}$\end{adjustbox}
\end{caja}
\end{columns}
\end{frame}

% ---------------------------------------------------------------------
\begin{frame}{{\normalsize Corriente total prevista del RLC en paralelo con la bobina}}
\begin{columns}[c,onlytextwidth]
  \column{0.62\textwidth}
\centering
\begin{tikzpicture}
  \begin{semilogxaxis}[width=8.4cm, height=5.7cm, xmin=2000, xmax=15000, ymin=0, ymax=17.334,
      xlabel={Frecuencia f (Hz)}, ylabel={Corriente total IT (mA)},
      label style={font=\scriptsize, text=iubnavy}, tick label style={font=\tiny, text=iubnavy},
      axis line style={iubnavy}, grid=major, grid style={iubnavy!12},
      legend style={font=\tiny, fill=white}, legend pos=north east]
    \addplot[line width=1.4pt, iubblue] coordinates {(2000,12.732) (2054.6,12.276) (2110.8,11.828) (2168.5,11.387) (2227.7,10.954) (2288.6,10.527) (2351.1,10.107) (2415.3,9.6931) (2481.3,9.2857) (2549.1,8.8845) (2618.8,8.4891) (2690.3,8.0993) (2763.8,7.715) (2839.3,7.3358) (2916.9,6.9617) (2996.6,6.5923) (3078.5,6.2276) (3162.6,5.8674) (3249,5.5114) (3337.7,5.1596) (3428.9,4.8118) (3522.6,4.468) (3618.9,4.1281) (3717.7,3.7922) (3819.3,3.4603) (3923.7,3.1328) (4030.9,2.8101) (4141,2.4931) (4254.1,2.1834) (4370.4,1.8834) (4489.8,1.5977) (4612.4,1.3344) (4738.5,1.1088) (4867.9,0.9474) (5000.9,0.88572) (5137.6,0.94341) (5277.9,1.1021) (5422.1,1.3264) (5570.3,1.5894) (5722.5,1.8754) (5878.8,2.1761) (6039.4,2.4869) (6204.4,2.8053) (6374,3.1296) (6548.1,3.4591) (6727,3.7932) (6910.8,4.1317) (7099.6,4.4744) (7293.6,4.8214) (7492.9,5.1727) (7697.6,5.5285) (7907.9,5.8887) (8124,6.2537) (8345.9,6.6236) (8574,6.9987) (8808.2,7.3791) (9048.9,7.7651) (9296.1,8.1569) (9550.1,8.5549) (9811,8.9592) (10079,9.3702) (10354,9.788) (10637,10.213) (10928,10.646) (11227,11.086) (11533,11.535) (11848,11.992) (12172,12.457) (12505,12.932) (12846,13.417) (13197,13.911) (13558,14.415) (13928,14.93) (14309,15.455) (14700,15.992) (15000,16.403)};
    \addlegendentry{IT prevista}
    \addplot[line width=1pt, dashed, iubnavy!70] coordinates {(4976,0) (4976,17.334)};
    \addlegendentry{fr \ensuremath{\approx} 4.98 kHz}
  \end{semilogxaxis}
\end{tikzpicture}
  \column{0.35\textwidth}
\begin{caja}{iubblue}{Lectura de la gráfica}
\footnotesize Es la corriente de circulación del tanque, que no pasa por la fuente.
\end{caja}
\end{columns}
\end{frame}

% ---------------------------------------------------------------------
\begin{frame}{RLC en paralelo del kit con VS = 2.00 V eficaces}
\centering\tiny
\renewcommand{\arraystretch}{1.35}
\rowcolors{2}{iubnavy!6}{white}
\begin{adjustbox}{max width=\linewidth, max totalheight=5.6cm}
\begin{tabularx}{\textwidth}{>{\raggedright\arraybackslash}X>{\raggedright\arraybackslash}X>{\raggedright\arraybackslash}X>{\raggedright\arraybackslash}X>{\raggedright\arraybackslash}X>{\raggedright\arraybackslash}X>{\raggedright\arraybackslash}X}
  \rowcolor{iubnavy}
  \hd{Magnitud} & \hd{3 kHz calc.} & \hd{3 kHz med.} & \hd{fr calc. (4.98 kHz)} & \hd{fr med. (4.96 kHz)} & \hd{8 kHz calc.} & \hd{8 kHz med.}\\
  IR (mA) & 0.43 & 0.43 & 0.43 & 0.43 & 0.43 & 0.43\\
  Ibob (mA) & 10.12 & 10.1 & 6.13 & 6.10 & 3.82 & 3.81\\
  IC (mA) & 3.69 & 3.70 & 6.12 & 6.08 & 9.83 & 9.82\\
  IT (mA) & 6.58 & 6.57 & 0.885 & 0.89 & 6.05 & 6.04\\
  Ángulo de IT (°) & \ensuremath{-}75.2 & — & \ensuremath{\approx} 0 & — & +84.3 & —\\
  Carácter & inductivo & — & resistivo & — & capacitivo & —\\
\end{tabularx}
\end{adjustbox}
\vspace{2mm}
{\scriptsize RLC en paralelo del kit con VS = 2.00 V eficaces: cálculo con la bobina real y medición (lecturas de ejemplo)\par}
\end{frame}

% =====================================================================
\bloque{08}{Síntesis, cierre actitudinal y bibliografía}{15 min}{iuborange}
% =====================================================================

% ---------------------------------------------------------------------
\begin{frame}{Síntesis de la sesión}
\begin{columns}[T,onlytextwidth]
  \column{0.32\textwidth}
\begin{cardB}{Admitancia e impedancia del circuito RLC}
\scriptsize Y = G + j(BC \ensuremath{-} BL), carácter inductivo o capacitivo según la frecuencia
\end{cardB}
  \column{0.32\textwidth}
\begin{cardO}{Ley de corrientes de Kirchhoff}
\scriptsize IT = IR + j(IC \ensuremath{-} IL): corrientes reactivas en oposición y mayores que la total
\end{cardO}
  \column{0.32\textwidth}
\begin{cardG}{Resonancia en paralelo}
\scriptsize fr, Z = R, Q = R/XL, ancho de banda y curva de impedancia
\end{cardG}
\end{columns}
\vspace{2mm}
\begin{columns}[T,onlytextwidth]
  \column{0.32\textwidth}
\begin{cardB}{El circuito tanque con bobina real}
\scriptsize Equivalente RLC en paralelo, Rp(eq) = RW(Q² + 1) y desplazamiento de la resonancia
\end{cardB}
  \column{0.32\textwidth}
\begin{cardO}{Corrección del factor de potencia}
\scriptsize Compensar la corriente reactiva de una carga inductiva y dimensionar el condensador
\end{cardO}
  \column{0.32\textwidth}
\begin{cardG}{Medición de un circuito RLC en paralelo}
\scriptsize Con el kit: 4.7 k\ensuremath{\Omega}, bobina de 10.4 mH y condensador de 97.8 nF a 3 kHz, fr y 8 kHz
\end{cardG}
\end{columns}
\end{frame}

% ---------------------------------------------------------------------
\begin{frame}{Reflexión para llevar}
\centering
\begin{tikzpicture}
  \node[fill=iubnavy, text=white, rounded corners=6pt, text width=11.5cm, inner sep=12pt, align=center, font=\normalsize] (c) at (0,0) {\itshape «Instalar un banco de condensadores sin calcular antes la corriente que circulará por él, o sin sus resistencias de descarga, pone en riesgo a quien lo intervenga después: el trabajo técnico responsable piensa también en el compañero del siguiente turno.»};
  \node[fill=iubyellow, text=iubnavy, rounded corners=3pt, font=\scriptsize\bfseries, inner sep=4pt, anchor=south west] at ([yshift=-4pt]c.north west) {Componente actitudinal};
\end{tikzpicture}
\end{frame}

% ---------------------------------------------------------------------
\begin{frame}{Referencias bibliográficas}
\small
\begin{itemize}\setlength{\itemsep}{8pt}
  \item[{$[1]$}] T. L. Floyd, Principios de circuitos eléctricos, 8.\textordfeminine{} ed. México: Pearson Educación, 2007.
  \item[{$[2]$}] R. L. Boylestad, Introducción al análisis de circuitos, 10.\textordfeminine{} ed. México: Pearson Educación, 2004.
  \item[{$[3]$}] M. F. P. Deorsola y P. Morcelle del Valle, Circuitos eléctricos. Parte 1, 1.\textordfeminine{} ed. La Plata: Editorial de la Universidad de La Plata, 2017.
\end{itemize}
\end{frame}

% ---------------------------------------------------------------------
%  CIERRE
% ---------------------------------------------------------------------
\begin{frame}[plain]
\begin{tikzpicture}[remember picture, overlay]
  \fill[iubnavy] (current page.north west) rectangle (current page.south east);
  \node[anchor=north west, text=iubyellow, font=\bfseries\fontsize{40}{44}\selectfont]
    at ([xshift=1.2cm,yshift=-1.6cm]current page.north west) {\textexclamdown{}Gracias!};
  \node[anchor=north west, text=white, font=\normalsize, text width=14cm, align=left]
    at ([xshift=1.2cm,yshift=-3.4cm]current page.north west)
    {IEL04 \textperiodcentered{} Circuitos Eléctricos I};
  \node[anchor=north west, fill=iubblue, text=white, rounded corners=1.5mm,
        font=\small\bfseries, inner sep=2.2mm, text width=11.5cm]
    at ([xshift=1.2cm,yshift=-4.4cm]current page.north west)
    {Próxima sesión \textperiodcentered{} Semana 13: Cálculo y medición de parámetros en circuitos RLC serie y paralelo: unidades, simbología y nomenclatura};
  \node[anchor=south west, text=iubyellow, font=\small, align=left]
    at ([xshift=1.2cm,yshift=0.9cm]current page.south west)
    {Ing. Sergio Martinez Castilla \quad · \quad smartinezc@unibarranquilla.edu.co};
  \fill[iubblue]   ([xshift=1.2cm,yshift=0.55cm]current page.south west) rectangle ++(1.6cm,0.15cm);
  \fill[iuborange] ([xshift=2.9cm,yshift=0.55cm]current page.south west) rectangle ++(1.6cm,0.15cm);
  \fill[iubgreen]  ([xshift=4.6cm,yshift=0.55cm]current page.south west) rectangle ++(1.6cm,0.15cm);
\end{tikzpicture}
\end{frame}

\end{document}
